
The computational complexity of the algorithm is $O(md_1d_2 + d_2^3)$ for stage 1, while stage 2 requires additional $O(nd_1d_2 + d_1^3)$ computations. This is small compared to KIV \citep{Singh2019}, which takes $O(m^3)$ and $O(n^3)$, respectively. We can further speed up the learning by using mini-batch training as shown in Algorithm~\ref{algo}. Here, $\hat{\vec{V}}^{(m_b)}$ and $\hat{\vec{u}}^{(n_b)}$ are the functions given by \eqref{eq:stage1-sol} and \eqref{eq:stage2-sol} calculated using mini-batches of data. Similarly, $\mathcal{L}_1^{(m_b)}$ and $\mathcal{L}_1^{(n_b)}$ are the stage 1 and 2 losses for the mini-batches. We recommend setting the batch size large enough so that $\hat{\vec{V}}^{(m_b)}, \hat{\vec{u}}^{(n_b)}$ do not diverge from $\hat{\vec{V}}^{(m)}, \hat{\vec{u}}^{(n)}$ computed on the entire dataset. Furthermore, we observe that setting $T_1 > T_2$, i.e. updating $\theta_Z$ more frequently than $\theta_X$, stabilizes the learning process.

 \begin{algorithm}
 \caption{Deep Feature Instrumental Variable Regression}
 \begin{algorithmic}[1]  \label{algo}
 \renewcommand{\algorithmicrequire}{\textbf{Input:}}
 \renewcommand{\algorithmicensure}{\textbf{Output:}}
 \REQUIRE Stage 1 data $(x_i, z_i)$, Stage 2 data $(\tilde y_i \tilde z_i)$, Regularization parameters $(\lambda_1, \lambda_2)$. Initial values $\hat{\theta}_X,\hat{\theta}_Z$. Mini-batch size $(m_b, n_b)$. Number of updates in each stage $(T_1, T_2)$.
 \ENSURE Estimated structural function $\hat{f}_\mathrm{struct}(x)$
 \REPEAT
   \STATE Sample $m_b$ stage 1 data $(x^{(b)}_i, z^{(b)}_i)$ and  $n_b$ stage 2 data $(\tilde y^{(b)}_i, \tilde z^{(b)}_i)$.
  \FOR{$t=1$ to $T_1$}
    \STATE Return function $\hat{\vec{V}}^{(m_b)}(\hat{\theta}_X,\theta_Z)$ in \eqref{eq:stage1-sol}  using $(x^{(b)}_i, z^{(b)}_i)$
    \STATE Update $\hat{\theta}_Z \leftarrow \hat{\theta}_Z - \alpha \nabla_{\theta_Z} \mathcal{L}^{(m_b)}_1(\hat{\vec{V}}^{(m_b)}(\hat{\theta}_X,\theta_Z), \theta_Z)|_{\theta_Z=\hat{\theta}_Z}$ \hspace{5mm}\textit{$\backslash\backslash$ Stage 1 learning}
  \ENDFOR
  \FOR{$t=1$ to $T_2$}
    \STATE Return function $\hat{\vec{u}}^{(n_b)}(\theta_X,\hat{\theta}_Z)$ in  \eqref{eq:stage2-sol} using $ (\tilde y^{(b)}_i, \tilde z^{(b)}_i)$ and function $\hat{\vec{V}}^{(m_b)}(\theta_X,\hat\theta_Z)$\!
    \STATE Update $\hat{\theta}_X \leftarrow \hat{\theta}_X -\alpha\nabla_{\theta_X} \mathcal{L}^{(n_b)}_2(\hat{\vec{u}}^{(n_b)}(\theta_X,\hat{\theta}_Z), \theta_X)|_{\theta_X=\hat{\theta}_X}$ \hspace{5mm}\textit{$\backslash\backslash$ Stage 2 learning}
  \ENDFOR
  \UNTIL{\textbf{convergence}}
  \STATE Compute $\hat{\vec{u}}^{(n)}:=\hat{\vec{u}}^{(n)}(\hat{\theta}_X,\hat{\theta}_Z)$ from \eqref{eq:stage2-sol} using entire dataset.
 \RETURN $\hat{f}_\mathrm{struct}(x) = (\hat{\vec{u}}^{(n)})^\top \vec{\psi}_{\hat{\theta}_X}(x)$
 \end{algorithmic} 
 \end{algorithm}


In Appendix~\ref{sec:consistency}, we provide regularity conditions under which the function learned by DFIV converges to the true structural function in probability. The derivation is based on Rademacher complexity bounds \citep{FoundationOfML}.%, which is applicable not only to neural networks but also to various machine learning techniques.